\documentclass[text05.tex]{subfiles}
\begin{document}
\newpage
\section{赤外線で信号を送る}
\subsection{赤外線ってなんだろう？}

\begin{figure}[H]
\centering
\includesvg[width=0.8\linewidth]{../../asset/image/electromagnetic_wave.svg}
    \caption{\ruby{電磁}{でん|じ}スペクトル図}

\end{figure}

\ruby{物質}{ぶっ|しつ}や生物は電磁波を発しています。X線、\ruby{紫外線}{し|がい|せん}、光、マイクロ波などは電磁波の一種です。電磁波は波であり、波の長さで種類をわけることができます。わたしたちは\ruby{可視光線}{か|し|こう|せん}と\ruby{呼}{よ}ばれる光が目に入ることで、色などを目で感知することができます。マイクロ波は電子レンジに使われています。マイクロ波によって物質を\ruby{振動}{しん|どう}させることで、物質の温度をあげています。\\

赤外線は、人の目には見えない光、つまり電磁波の\ruby{仲間}{なか|ま}です。温度のある物体は、目に見えない電磁波を\ruby{周囲}{しゅう|い}に出しており、その中には赤外線も\ruby{含}{ふく}まれています。多くの場合、物体の温度が高いほど、赤外線を強く出します。Raspberry Piも、動作中に電気の一部が熱に変わって温かくなり、赤外線を出しています。Raspberry Piに\ruby{触}{ふ}れなくても手を近づけると温かく感じることがあるのは、赤外線が手に\ruby{吸収}{きゅう|しゅう}されることや、温められた空気から熱が伝わることがあるためです。手が赤外線そのものを感じているわけではなく、皮ふが温まったことを感じています。\\\\

赤外線は、カメラを通すと人の目にも見えるようになります。カメラのセンサーが赤外線をとらえるからです。しかし、近年のデジタルカメラやスマートフォンのカメラには赤外線カットフィルターがついているものが多く、赤外線を見ることができません。配布したWebカメラには赤外線LEDの光が映ります。

\begin{figure}[H]
    \centering
    \includesvg[scale=0.8]{../../asset/image/infrared_rays_led.svg}
        \caption{Webカメラで\ruby{撮影}{さつ|えい}した赤外線LED}

\end{figure}

赤外線を使って信号を送ったり、受け取ったりすることもできます。例えば赤外線が送られているときは電気をつける、赤外線が送られていないときは電気を消す、などとすることができます。しかし赤外線が送られているかどうかを判定するだけではつける、消すの２つの動作しかできません。ロボットのように歩く、止まる、右に曲がる、左に曲がるなど多くの動きを\ruby{制御}{せい|ぎょ}したいときは、赤外線のONとOFFの組み合わせで信号を作って送ります。

\begin{figure}[H]
\centering
\includesvg[scale=0.8]{../../asset/image/text05-img037.svg}
\caption{赤外線で信号を送る方法}
\label{ir}
\end{figure}

図\ref{ir}を見てみましょう。電球が点いているときは赤外線を送り、電球が消えているときは赤外線を送っていないという意味です。赤外線リモコンでは、赤外線を高速で点滅させた短い光のまとまりを、数百から数千マイクロ秒送ったり止めたりします。つけたり消したりを組み合わせて、いろいろなパターンを作ることができます。そのパターンによって動作を決めています。送る、送らない、送らない、送らない、送る、送る、送らないと信号が来たときは”歩く”動作をします。
\end{document}
